% MAS1004 Assignment 2: long report.
%
% You fill this in with your agent, in your repository. AGENTS.md tells the
% agent how. Red text in the PDF is what still has to be replaced.
%
% - Pictures go in figures/. Copy them from results/, which is not in git.
% - Outputs go in outputs/: copy experiments.csv, LOG.md and the outputs of
%   check_submission.py and curate.py count here. \outputfile{outputs/...}
%   prints a file exactly as it is. Every number in this report has to come
%   from such an output.
%
% To make the PDF: compress this folder into a zip, and on Overleaf choose
% New Project, Upload Project. Set Menu, Compiler to XeLaTeX, and Recompile.

\documentclass[11pt]{article}
\usepackage[ko]{mas1004}

\reporttitle{Assignment 2: Long report}
\studentname{\fillin{이름}}
\studentid{\fillin{학번}}

\begin{document}
\makeheader

리더보드 닉네임: \fillin{닉네임}.
마지막 릴리스: Track A \texttt{track-a-2} (``no horizontal flip'', 레시피 \texttt{no\_flip}),
Track B \texttt{track-b-1} (``MobileNetV3-L 320, balanced cities'', 모델 \texttt{b1}).

학습은 모두 Google Colab의 T4 GPU에서 했다. 내 노트북(Apple M5)은 Mac GPU(mps)에서
학습이 오류로 멈췄고(원래 레시피로도 같은 오류), CPU로는 Track A 한 번에 약 193분이
걸려서 쓰지 않았다. 무료 Colab의 GPU 사용량이 하루 만에 끝나서, 뒤쪽 실행
(\texttt{no\_flip}, \texttt{b1})은 다른 Google 계정의 Colab에서 돌렸다. 같은 저장소의 같은 커밋을
가져다 쓰므로 실행 조건은 같다.

% ---------------------------------------------------------------------------
\section*{Track A}

\begin{center}\footnotesize
\begin{tabular}{lp{0.5\linewidth}r}
\toprule
tag & what changed from the run before & validation \\
\midrule
baseline (seed 0) & starter recipe 그대로 & 0.4981 \\
baseline (seed 1) & seed만 바꿈 & 0.4924 \\
no\_flip (seed 0) & \texttt{train\_transform}에서 \texttt{T.RandomHorizontalFlip()} 삭제 & 0.4933 \\
\bottomrule
\end{tabular}
\end{center}

\paragraph{baseline.} starter recipe를 seed 0과 1로 두 번 학습했다. 바뀐 것은 seed뿐인데
점수가 0.4981과 0.4924로 0.0057 차이가 났다. 그래서 이보다 작은 차이는 운일 수 있다고 보고
다음 실행들을 비교했다.

도시별로 보면 두 seed 모두 부산(0.375, 0.385)과 오사카(0.358, 0.375)가 가장 낮았다.
seed 0의 confusion matrix(\texttt{results/baseline\_seed0/result.json})를 보면 부산 검증 사진
1,390장 중 521장을 맞혔고 444장을 서울로 틀렸다. 오사카는 2,128장 중 762장을 맞혔고
844장을 도쿄로 틀려서, 맞힌 것보다 도쿄로 틀린 것이 더 많았다. 학습 사진 수도
부산 5,191장, 오사카 7,200장으로 서울(11,274장)과 도쿄(13,347장)보다 훨씬 적다.

\paragraph{no\_flip.} 이유: ``글자가 중요한 단서일 것 같은데, 한글, 일본어, 한자는 뒤집으면 다른
글자가 되니까 구별이 잘 안 될 것 같아.'' 뒤집기를 빼면 차가 다니는 방향(도쿄·오사카는 좌측통행)도
그대로 남는다.

결과는 0.4933으로 baseline 두 번(0.4924, 0.4981) 사이에 있어서, 좋아졌다고도 나빠졌다고도
할 수 없다. 서울은 0.586으로 조금 올랐고 오사카는 0.342로 조금 내렸다.
내 해석: ``모델이 글자보다 건물 모양, 색감, 도로 같은 다른 단서를 더 많이 쓰고 있을 수 있다.
그리고 같은 나라 끼리라서 헷갈리는 듯.''

\paragraph{최종 레시피.} \texttt{no\_flip}을 냈다(selection.csv 없음, 모든 학습 사진 사용).
검증 점수는 baseline과 같은 수준이지만, baseline은 starter recipe 그 자체라 이길 수 없고,
리더보드와 최종 채점 사진(길거리 영상, 휴대폰 사진)에는 간판 글자와 차선 방향이 실제로 담겨 있다.
리더보드 점수: \fillin{다음 날 아침 리더보드에 나오는 track-a-2 점수. starter recipe(track-a-1)의 점수도 같이}.

% ---------------------------------------------------------------------------
\section*{Decisions you made from photos}

\picturefile{figures/mistakes_baseline_seed0_busan_00.png}
\picturefile{figures/mistakes_baseline_seed0_osaka_00.png}

baseline (seed 0) 모델이 가장 확신하면서 틀린 검증 사진들(\texttt{curate.py mistakes})을
도시별로 한 장씩 보았다. 위는 부산과 오사카 사진이다. 사진 아래의 \texttt{-> seoul 0.99}는
``서울이라고 99\% 확신했는데 틀렸다''는 뜻이다.

내가 본 공통점: 글자가 안 보이는 사진, 전통 건물, 고속도로를 잘 구별하지 못하는 듯하다.
이것은 \texttt{no\_flip}의 결과와도 이어진다. 글자가 없는 사진이 많다면, 글자를 지키려고 뒤집기를
뺀 것이 큰 차이를 만들지 못한 것도 설명이 된다.

이 관찰로 규칙을 정해 학습 사진을 바꾸고 측정하는 단계까지는 가지 못했다.
첫날 Colab 세션이 사진 링크 입력을 기다리다 끊기고, 다음 날은 무료 GPU 사용량이 바닥나서
남은 GPU 시간을 \texttt{no\_flip}과 Track B에 썼기 때문이다. \fillin{한 주가 더 있다면 이 관찰로
무엇을 해 보고 싶은지, 본인 말로}

에이전트에게 한 요청:
\begin{itemize}
\item 에이전트가 Colab에서 만든 baseline의 틀린 사진 시트(\texttt{curate.py mistakes})를 직접 보았다.
\item ``뒤집어보자. 왜냐면 내가 생각했을땐 글자가 중요한 단서일 것 같은데, 한글, 일본어,
  한자는 뒤집으면 다른 글자가 되니까 구별이 잘 안 될 것 같아.'' $\rightarrow$ 레시피에서
  \texttt{RandomHorizontalFlip} 삭제(\texttt{no\_flip}).
\item Track B는 에이전트가 제안한 설정(MobileNetV3-Large, 320, 6 epochs, 도시별로 고르게 뽑기)에
  동의해서 돌렸다.
\item Track A 최종 레시피는 에이전트에게 추천을 받아 \texttt{no\_flip}으로 정했다.
\end{itemize}

% ---------------------------------------------------------------------------
\section*{Track B}

\begin{center}\footnotesize
\begin{tabular}{lp{0.5\linewidth}r}
\toprule
tag & what & validation \\
\midrule
b1 & MobileNetV3-Large (IMAGENET1K\_V2), 320 $\times$ 320, 6 epochs, 도시별로 고르게 뽑기(\texttt{--balance}), starter 학습 사진만 & 0.554 \\
\bottomrule
\end{tabular}
\end{center}

\texttt{src/train\_b.py}로 한 번 학습했다. 4.20 million parameters, 0.44 billion multiply-adds로
한도(12 million, 2.0 billion) 안에 든다. 학습에 약 59분(3,540초) 걸렸다.

이유: ``도시별로 고르게 뽑자고 한 이유는, 한 도시만 많으면 그 도시에 편향될 줄 알았다.''

검증 점수는 0.554로 Track A baseline(0.4981)보다 높다. 도시별로는 부산 0.375 $\rightarrow$ 0.493,
오사카 0.358 $\rightarrow$ 0.491로 사진이 적은 도시가 많이 올랐고, 사진이 가장 많은 도쿄는
0.641 $\rightarrow$ 0.604로 내려갔다(Track A baseline seed 0과 비교).
내 해석: ``지금 생각해보면 그 도시에 대한 정확성은 좀 오르지 않았을까... 그리고 그래도 글자의
영향이 있었기 때문에 높게 나온 듯.''

다만 \texttt{b1}은 모델, 이미지 크기, 도시별 균형, 학습량을 한꺼번에 바꿨기 때문에, 이 중 무엇이
점수를 올렸는지는 이 실행만으로는 나눠 말할 수 없다.

추가 사진은 모으지 않았다.

검증 점수와 리더보드 점수: \fillin{track-b-1의 리더보드 점수, 그리고 두 점수가 다를 수 있는 이유에 대한 본인 생각}.

\subsection*{Your submission, as the grader sees it}
\outputfile{outputs/check_submission.txt}

% ---------------------------------------------------------------------------
\section*{What made the difference}

가장 큰 차이는 Track B의 \texttt{b1}이었다: 0.554로 Track A baseline보다 약 0.056 높다.
다만 여러 가지를 함께 바꾼 결과다. Track A의 \texttt{no\_flip}은 0.4933으로, seed만 바꿨을 때의
차이(0.0057) 안에 있어서 효과가 있었다고 할 수 없다.

\fillin{한 주가 더 있다면 무엇을 해 보고 싶은지, 본인 말로}

% ---------------------------------------------------------------------------
\section*{All runs}
\outputfile{outputs/experiments.csv}

\section*{Lab notebook}
\outputfile{outputs/LOG.md}

\end{document}
